% concatenated from main.tex
\documentclass{amsart}

\usepackage{amsmath,amssymb,amsthm}
\usepackage{fullpage}

\newtheorem{theorem}{Theorem}
\newtheorem{lemma}{Lemma}
\newtheorem{proposition}{Proposition}
\newtheorem{corollary}{Corollary}

\title{(Claw, $\vec C_3$)-free Digraphs with Unbounded Dichromatic Number}

\author{Guillaume Aubian}
\address{CRED, Université Paris-Panthéon-Assas}
% \email{guillaume.aubian@...}

\author{Luis Kuffner}
\address{Université Paris Cité, CNRS, IRIF, F-75013 Paris, France}
% \email{luis.kuffner@...}

\date{November 2025}

\begin{document}
\maketitle


\begin{abstract}
We construct orientations of rook graphs (whose underlying graphs are claw-free) that contain no $\vec C_3$ but have unbounded dichromatic number. This disproves a conjecture of Aboulker, Charbit and Naserasr and improves a result of Carbonero, Koerts, Moore and Spirkl.
\end{abstract}


\section{Introduction and notation}



A hero is a tournament $H$ such that tournaments not containing $H$ have bounded dichromatic number. In a seminal paper \cite{BCCFLSST13}, Berger, Choromanski, Chudnovski, Fox, Loebl, Scott, Seymour and Thomassé obtained a structure theorem for heroes. Harutyunyan, Le, Newman and Thomassé \cite{HLNT17} later proved that oriented graphs with bounded independence number (tournaments being the special case where the independence number equals one) not containing a hero have bounded dichromatic number. Chudnovsky, Scott and Seymour \cite{CSS19} subsequently showed that for any transitive tournament $H$ and oriented star $S$, oriented graphs containing neither $S$ nor $H$ as induced subdigraphs have bounded dichromatic number. Aboulker, Charbit and Naserasr \cite{ACN21} thus wondered whether (and actually conjectured that), for any hero $H$ and any oriented star forest $S$, oriented graphs containing neither $S$ nor $H$ as induced subdigraphs have bounded dichromatic number.

\smallskip{}

Aboulker, Aubian and Charbit \cite{AAC24} provided a counterexample when $S = K_1 + \vec K_2$. Carbonero, Koerts, Moore and Spirkl \cite{CKMS25} did so when $H = \Delta(1,1,C_3)$ and $S$ is any orientation of $K_{1,4}$, and also when $H = \vec C_3$ and $S$ is any orientation of $K_{1,5}$. We improve these two latter results with a counterexample when $H = \vec C_3$ and $S$ is any orientation of $K_{1,3}$, i.e. the claw.

\bigskip{}

Definitions in this paper are standard and follow from Bondy and Murty's book \cite{BM08}. In this paper, we study oriented graphs, i.e. directed graphs with no pair of opposite arcs. Given two integers $n,i \ge 0$, we write $n_i$ for the $i$th least significant bit of the binary expansion of $n$, starting from $i=0$.

\section{Construction and proof}

Let $N$ be a positive integer. The \emph{$N \times N$ rook graph} $R_N$ is the graph with vertex set $[1,N]^2$ and with an edge between two vertices whenever they share their first coordinate or they share their second coordinate. Intuitively, it is the graph describing how a rook moves on a $N \times N$ chessboard. It is the line graph of $K_{N,N}$, and therefore is claw-free. We orient it into a digraph $D_N$ as follows. For any vertices $(a,b),\ (a,d),\ (c,b)$, we orient the edge from $(a,b)$ to $(a,d)$ if and only if $b_i = a_i$, where $i = \min\{ j \mid b_j \ne d_j \}$, and we orient the edge from $(a,b)$ to $(c,b)$ if and only if $b_i \ne a_i$, where $i = \min\{ j \mid a_j \ne c_j \}$.

\begin{proposition}
The digraph $D_N$ contains no $\vec C_3$.
\end{proposition}
\begin{proof}
Given three vertices $(a,b)$, $(a,c)$ and $(a,d)$ on the same row, let $i = \min\{j \mid \{b_j, c_j, d_j\} = \{0,1\}\}$. Without loss of generality let us assume that $b_i \ne c_i = d_i$. Edges $(a,b)(a,c)$ and $(a,b)(a,d)$ are oriented only depending on whether $a_i = b_i$ and thus have the same orientation. Thus, $(a,b)$, $(a,c)$ and $(a,d)$ can not form a $\vec C_3$. The same argument applies to any column. Since every triangle in a rook graph is contained entirely in a single row or in a single column, $D_N$ contains no $\vec C_3$.
\end{proof}

\begin{lemma}
When $|a - c| = |b - d|$, vertices $(a,b),\ (a,d),\ (c,d),\ (c,b)$ induce a directed $4$-cycle.
\end{lemma}
\begin{proof}
Let $i=\min \{j \mid |a - c|_{j}  = 1\}$. Then $i = \min\{j \mid a_j \ne c_j\} = \min\{j \mid b_j \ne d_j\}$.
Up to permuting $a$ and $c$ and/or permuting $b$ and $d$, we may assume $a_i = b_i = 1$ and $c_i = d_i = 0$.
Since $b_i = a_i$, we have $(a,b) \to (a,d)$. Since $d_i \ne a_i$, we have $(a,d) \to (c,d)$. Since $d_i = c_i$, we have $(c,d) \to (c,b)$. Since $b_i \ne c_i$, we have $(c,b) \to (a,b)$.
Thus the four vertices form the directed cycle $(a,b) \to (a,d) \to (c,d) \to (c,b) \to (a,b)$.
\end{proof}

Note that it actually suffices that $d_2(a,c) = d_2(b,d)$ where $d_2$ is the dyadic distance, but this stronger condition suffices for the sake of our result.

\begin{theorem}
The family $(D_N)_{N\ge 1}$ has unbounded dichromatic number.
\end{theorem}

\begin{proof}
Fix $k$. The Gallai--Witt theorem  \cite{W51} (also known as the multidimensional van der Waerden theorem) implies that there exists $N$ such that any $k$-colouring of $[1,N]^2$ contains a monochromatic axis-parallel square, that is, four points of the form $(x,y),\ (x+t,y),\ (x,y+t),\ (x+t,y+t)$ for some integers $x,y,t>0$. Thus, for $N$ large enough, any colouring of $V(D_N)$ with $k$ colours contains a monochromatic square as above. By the previous lemma, the four vertices of such a square induce a directed $4$-cycle. Hence, some colour class does not induce an acyclic subdigraph. Since $k$ is arbitrary, the dichromatic number of $D_N$ is unbounded.
\end{proof}


\bibliographystyle{amsplain}
\bibliography{main}

\end{document}
